% Options for packages loaded elsewhere
\PassOptionsToPackage{unicode}{hyperref}
\PassOptionsToPackage{hyphens}{url}
\documentclass[
  ignorenonframetext,
]{beamer}
\newif\ifbibliography
\usepackage{pgfpages}
\setbeamertemplate{caption}[numbered]
\setbeamertemplate{caption label separator}{: }
\setbeamercolor{caption name}{fg=normal text.fg}
\beamertemplatenavigationsymbolsempty
% remove section numbering
\setbeamertemplate{part page}{
  \centering
  \begin{beamercolorbox}[sep=16pt,center]{part title}
    \usebeamerfont{part title}\insertpart\par
  \end{beamercolorbox}
}
\setbeamertemplate{section page}{
  \centering
  \begin{beamercolorbox}[sep=12pt,center]{section title}
    \usebeamerfont{section title}\insertsection\par
  \end{beamercolorbox}
}
\setbeamertemplate{subsection page}{
  \centering
  \begin{beamercolorbox}[sep=8pt,center]{subsection title}
    \usebeamerfont{subsection title}\insertsubsection\par
  \end{beamercolorbox}
}
% Prevent slide breaks in the middle of a paragraph
\widowpenalties 1 10000
\raggedbottom
\AtBeginPart{
  \frame{\partpage}
}
\AtBeginSection{
  \ifbibliography
  \else
    \frame{\sectionpage}
  \fi
}
\AtBeginSubsection{
  \frame{\subsectionpage}
}
\usepackage{iftex}
\ifPDFTeX
  \usepackage[T1]{fontenc}
  \usepackage[utf8]{inputenc}
  \usepackage{textcomp} % provide euro and other symbols
\else % if luatex or xetex
  \usepackage{unicode-math} % this also loads fontspec
  \defaultfontfeatures{Scale=MatchLowercase}
  \defaultfontfeatures[\rmfamily]{Ligatures=TeX,Scale=1}
\fi
\usepackage{lmodern}
\ifPDFTeX\else
  % xetex/luatex font selection
\fi
% Use upquote if available, for straight quotes in verbatim environments
\IfFileExists{upquote.sty}{\usepackage{upquote}}{}
\IfFileExists{microtype.sty}{% use microtype if available
  \usepackage[]{microtype}
  \UseMicrotypeSet[protrusion]{basicmath} % disable protrusion for tt fonts
}{}
\makeatletter
\@ifundefined{KOMAClassName}{% if non-KOMA class
  \IfFileExists{parskip.sty}{%
    \usepackage{parskip}
  }{% else
    \setlength{\parindent}{0pt}
    \setlength{\parskip}{6pt plus 2pt minus 1pt}}
}{% if KOMA class
  \KOMAoptions{parskip=half}}
\makeatother
\setlength{\emergencystretch}{3em} % prevent overfull lines
\providecommand{\tightlist}{%
  \setlength{\itemsep}{0pt}\setlength{\parskip}{0pt}}
% Slide layout defaults for warning-clean scientific decks.
\usepackage{etoolbox}
\IfFileExists{xurl.sty}{\usepackage{xurl}}{}
\IfFileExists{seqsplit.sty}{\usepackage{seqsplit}}{\newcommand{\seqsplit}[1]{#1}}
\protected\def\breakseq#1{\seqsplit{#1}}
\protected\def\breaktt#1{\begingroup\ttfamily\seqsplit{#1}\endgroup}
\setlength{\emergencystretch}{6em}
\tolerance=5000
\hbadness=10000
\hfuzz=1pt
\setlength{\tabcolsep}{2pt}
\AtBeginEnvironment{longtable}{\tiny\renewcommand{\arraystretch}{0.86}\setlength{\tabcolsep}{1pt}}
\AtBeginEnvironment{tabular}{\tiny\renewcommand{\arraystretch}{0.86}\setlength{\tabcolsep}{1pt}}

% Natbib and cross-reference fallbacks — slides don't load natbib
% or cleveref, but combined-PDF manuscript prose may emit these
% commands. The fallback renders citations as a bracketed key list
% and cross-references as detokenized labels so slides don't fail on
% undefined control sequences or raw underscores. \providecommand is
% a no-op if packages load later.
\providecommand{\citep}[1]{[#1]}
\providecommand{\citet}[1]{#1}
\providecommand{\citealp}[1]{#1}
\providecommand{\citeauthor}[1]{#1}
\providecommand{\citeyear}[1]{#1}
\providecommand{\cref}[1]{\texttt{\detokenize{#1}}}
\providecommand{\Cref}[1]{\texttt{\detokenize{#1}}}

\newtheorem{proposition}{Proposition}
\newtheorem{hypothesis}{Hypothesis}
\usepackage{bookmark}
\IfFileExists{xurl.sty}{\usepackage{xurl}}{} % add URL line breaks if available
\urlstyle{same}
% fallback for those not using the hyperref driver hyperxmp:
\makeatletter
\@ifundefined{xmpquote}{\newcommand{\xmpquote}[1]{#1}}{}
\makeatother
\hypersetup{
  hidelinks,
  pdfcreator={LaTeX via pandoc}}

\author{\texorpdfstring{}{}}
\date{}

\begin{document}

\section{Relational Aesthetics: Coauthored Marks and Social
Form}\label{sec:relational_aesthetics}

\begin{frame}[allowframebreaks]{Relational Aesthetics: Coauthored Marks
and Social Form}
Relational aesthetics frames art as a production of human relations
rather than only objects for contemplation
\breakseq{[@bourriaud2002relational]}. Bishop's critique is equally important
because it warns against treating every engineered encounter as
automatically emancipatory or socially meaningful
\breakseq{[@bishop2004antagonism]}. PPPiP already had a relational structure:
the drawing mattered because it enacted the relationship. The couple
joint drawing method gives this point a nearby arts-therapy precedent by
treating shared drawing as an expression of connectedness and
individuality rather than a finished-image score \breakseq{[@snir2013joint]}.
DigiPPPiP makes that ontology explicit in digital form. The artifact is
a trace of mutual attention, negotiation, and improvisation; the canvas
is a social scene.

Digital-intimacy research sharpens the claim. Technologies for
strong-tie relationships often work through awareness, expressivity,
physicalness, gift giving, joint action, and memory rather than through
high-volume information transfer \breakseq{[@hassenzahl2012love;
@vetere2005mediating]}. DigiPPPiP sits closest to joint action and
memory: the partners do something together, and the resulting artifact
can remain available as a trace. That makes the practice more specific
than generic messaging and more cautious than therapeutic promise.

The first-principles relational test is whether the artifact still
points back to the other person as an agent. A canvas full of marks is
not relational if one partner dominates, if the other cannot perceive or
alter the shared field, or if the system substitutes automated novelty
for reciprocal response. Relational aesthetics therefore functions as a
constraint on design: the scene must leave room for mutual recognition,
refusal, repair, and surprise.

Creative arts therapy research offers a neurodynamic vocabulary for this
claim. Aesthetic engagement externalizes affect, supports co-regulation,
and creates intersubjective resonance \breakseq{[@vaisvaser2024neurodynamics]}.
Close-relationship co-regulation scholarship narrows the term: it should
mean bidirectional linkage across partners' emotional channels, not a
vague feeling of harmony \breakseq{[@butler2012coregulation]}.
Physiological-linkage reviews make the boundary even more conservative
because linkage can be helpful, neutral, or risky depending on task,
affective context, and overload \breakseq{[@timmons2015physiologicallinkage]}.
Family-level coregulation models add a useful multi-timescale warning:
dyadic regulation can be nested inside broader relational systems, so a
drawing trace alone cannot prove well-being
\breakseq{[@paley2022familycoregulation]}. Art-making studies also suggest
plausible stress and reflection constructs for future measurement
\breakseq{[@kaimal2016cortisol]}, while pandemic art-therapy discussions show
why remote creative practices need careful ethical framing
\breakseq{[@potash2020pandemics]}. Phygital art-therapy theory is useful here
because it treats digital and physical media as an integrated practice
ecology rather than as rival channels \breakseq{[@yoon2025phygital]}. DigiPPPiP
does not become therapy by declaration, but related mechanisms help
explain why shared drawing may be relationally meaningful: partners make
affect visible, modulate each other's timing, and jointly discover an
image neither would have made alone.

The important boundary is that DigiPPPiP remains human-human unless a
study or design explicitly examines AI mediation. Generative AI can
assist accessibility, prompts, or reflection, but it can also displace
the mutual recognition that makes partner drawing relational. AI
co-drawing belongs in a separate experimental branch, not as the default
form of the practice.
\end{frame}

\end{document}
